\begin{afterword}
\summaryhead

The post is a short mock film script. Generals learn that New York has been overrun by philosophical zombies, who behave exactly like people but are not conscious; video of the city before and after shows no difference. A scientist explains that the disease is an ``epiphenomenal virus'' with no effects, and no cure can touch it: placebos fail because they have effects.

David Chalmers, held in a glass cell as one of the first victims, says he feels fine and protests that the scenario distorts his views and cannot happen, but he is cut off before saying why. Daniel Dennett, who tells the police guarding Chalmers that he has no qualia, fights his way past them. The zombies answer the generals by saying that the humans are the ones with the virus; a colonel concludes that he feels nothing; the screen goes black and silent, and ``The movie continues exactly as before.'' A postscript shows the author attacked by zombie nurses at a convention.

\responsehead

The post is satire, and the question is whether its jokes are accurate about the view as its holders state it. For the most part they are.

The central joke is the definition itself. If consciousness had no effects, its loss would make no difference anyone could see: the city ``now'' looks like the city two weeks ago, the virus leaves no trace to study, and a cure would have to do nothing. Chalmers had written, in \textsc{The Conscious Mind}, that it was perhaps not surprising that phenomenal zombies had not been popular in Hollywood, given ``obvious problems with their depiction.'' The film makes those problems its joke.

The film is fair to Chalmers where it could easily not have been. The Chalmers character protests that the scenario ``can't actually happen,'' and that is Chalmers's view: \textsc{The Conscious Mind} holds that zombies are unlikely to be naturally possible, because in our world laws link physical states to experience. The film cuts him off before he can say so, but the next scene's ``Bridging Law Enforcement Agency'' names the laws.

Its one argument in miniature, that we know the virus exists ``The same way we know we're conscious,'' compresses the case made in ``Zombies! Zombies?'': on epiphenomenalism our knowledge of consciousness cannot come from its effects. Chalmers's reply, that having the experiences justifies the beliefs, is left out, as a joke may leave it out; the reply, and the author's answer to it, are in the earlier post.

The film adds no argument to the posts before it, and does not claim to. Its ending is its best line. The zombies' reply is the humans' own charge reversed, and when the picture and sound go, the film goes on unchanged, which is what the loss of epiphenomenal consciousness would amount to.

In short: a parody that turns the arguments of ``Zombies! Zombies?'' into jokes, is accurate to Chalmers's view where it states it, and ends with a line that puts the case against epiphenomenal consciousness more briefly than the argument did.
\end{afterword}
